\section{Síntesis y ampliación}

\subsection{Conclusiones principales}

Esta sección condensa en una lista de verificación los juicios sobre robótica de toda la nota. El lector puede usarla para evaluar una demo de robots o el discurso de una empresa: ¿resuelve el cerebro, el cerebelo, los datos, el cuerpo del robot, el tacto y la seguridad, o solo muestra una de estas piezas?

\begin{enumerate}
\item El primer cambio de paradigma de la robótica en la última década fue que el aprendizaje por refuerzo y Sim2Real transformaron la locomotion.
\item Los modelos grandes trajeron el segundo cambio de paradigma: dieron a los robots comprensión del lenguaje, sentido común y capacidad de planificación.
\item Los modelos base para robots siguen siendo, en lo esencial, modelos grandes multimodales a los que se agregan datos de acciones y fine-tuning; todavía no son del todo independientes.
\item Los datos de alta calidad son hoy el principal cuello de botella de la robótica; unos cientos de horas de datos no bastan para verificar el scaling.
\item Motion Transfer, de Gemini Robotics 1.5, intenta resolver la generalización entre distintos cuerpos de robot.
\item Synthetic Data es una dirección clave en la que, según Tan Jie (谭杰), hay que creer: los datos reales por sí solos no alcanzan.
\item El valor de los modelos del mundo está en predecir las consecuencias de las acciones, no solo en generar video.
\item Las manos diestras vuelven a dar importancia al tacto, porque la manipulation fina necesita información de contacto.
\item Los robots de propósito general pasarán por una etapa larga que va de la automatización fija a capacidades sobrehumanas.
\item La ruta ideal para la robótica de China y Estados Unidos es la complementariedad entre el paradigma de inteligencia y la cadena de suministro de hardware.
\end{enumerate}

\subsection{Preguntas abiertas}

Esta sección conserva preguntas abiertas porque muchas de las rutas técnicas discutidas en este episodio siguen cambiando con rapidez. Lo que más necesita la industria robótica no es una respuesta aislada, sino seguir de cerca qué variables mejoran de verdad: la escala de los datos, la transferencia entre cuerpos de robot, el hardware táctil, los modelos del mundo y la evaluación de seguridad.

\begin{itemize}
\item ¿Cuándo serán los modelos base para robots realmente independientes de los modelos grandes multimodales?
\item ¿Puede Motion Transfer funcionar de manera estable en un número suficiente de cuerpos de robot?
\item ¿Cómo puede Synthetic Data reducir la sim-to-real gap?
\item En los modelos del mundo, ¿importa más la predicción visual o la predicción de acciones y estados?
\item ¿Cuándo podrá el hardware táctil incorporarse a escala al ciclo de retroalimentación de datos de las manos diestras?
\item ¿Cuánto tiempo coexistirán los robots humanoides de propósito general y los specialist verticales?
\item ¿Cómo debe la safety de los robots adelantarse al rápido crecimiento de sus capacidades?
\end{itemize}

\subsection{Lecturas adicionales}

\begin{itemize}
\item Sim2Real Learning Agile Locomotion for Quadruped Robots: para entender cómo se transfiere el aprendizaje por refuerzo a robots cuadrúpedos reales.
\item Serie Robot Transformer: trabajos como RT-1, RT-2 y RT-X.
\item Gemini Robotics / Gemini Robotics 1.5: la línea de modelos para robots de Google DeepMind.
\item EP132, entrevista con Gao Jiyang (高继扬): Xinghaitu (星海图), Waymo/Momenta y la puesta en producción de la inteligencia corporizada.
\item EP134, panorama de datos: datos para robots, Recipe y fijación de precios de los datos.
\end{itemize}

\begin{importantbox}{Juicio final}
Lo más difícil en robótica no es «lograr que el modelo entienda una frase», sino lograr que un sistema físico actúe de forma estable en el mundo real. Los modelos grandes dieron a los robots un cerebro, el aprendizaje por refuerzo aportó parte del cerebelo, y los datos sintéticos y los modelos del mundo intentan suplir la experiencia; pero la implementación real todavía requiere datos de alta calidad, hardware confiable, tacto, seguridad y paciencia a largo plazo.
\end{importantbox}

\end{document}
